%% ma: L10-B14-0089
%% lop: 10
%% chuong: 5
%% bai: 14
%% dang: 10.14.12
%% loai: DS
%% muc: [NB, TH, TH, VD]
%% dap_an: "ĐSSS"
%% nguon: "chuyen-de-cac-so-dac-trung-cua-mau-so-lieu-khong-ghep-nhom-toan-10-chuong-trinh-moi.pdf (GV. Phan Nhật Linh) (P05), Bài 03 Dạng 2 Phần II Câu 6"
%% kien_thuc: "Lớp 10 Bài 13 (số trung bình, trung vị, tứ phân vị, mốt), Bài 14 (khoảng biến thiên, khoảng tứ phân vị, phương sai, độ lệch chuẩn, giá trị bất thường)"
%% trang_thai: chinh-thuc
%% ghi_chu: "Đáp án gốc ý a) và d) sai: lời giải gốc chép nhầm 450 thành 460, bỏ sót mốt 450 và ra phương sai 49375/2; đã sửa ý a) thành 'hai mốt 430 và 450' (Đ), ý d) là Sai. Tự duyệt theo yêu cầu giáo viên 22/10/2026."
\begin{ex}
Số lượng gạo bán ra tại một quầy hàng các tháng trong một năm được cho bởi bảng sau (đơn vị: yến):
\begin{center}
\begin{tabular}{|c|c|c|c|c|c|c|c|c|c|c|c|c|}\hline
Tháng & $1$ & $2$ & $3$ & $4$ & $5$ & $6$ & $7$ & $8$ & $9$ & $10$ & $11$ & $12$\\\hline
Số lượng & $410$ & $430$ & $450$ & $430$ & $525$ & $550$ & $950$ & $450$ & $800$ & $635$ & $760$ & $560$\\\hline
\end{tabular}
\end{center}
Xét tính đúng sai của các khẳng định sau:
\choiceTF{\True Mẫu số liệu có hai mốt là $430$ và $450$}{Trung vị của mẫu số liệu là $525$}{Khoảng tứ phân vị của mẫu số liệu là $257$}{Phương sai của mẫu số liệu là $s^2=\dfrac{49\,375}{2}$}
\loigiai{Sắp xếp: $410$; $430$; $430$; $450$; $450$; $525$; $550$; $560$; $635$; $760$; $800$; $950$.

a) Đúng: $430$ và $450$ cùng xuất hiện $2$ lần, nhiều nhất.

b) Sai: $M_e=\dfrac{525+550}2=537{,}5$.

c) Sai: $\Delta_Q=697{,}5-440=257{,}5$.

d) Sai: $\overline x=\dfrac{6950}{12}\approx579{,}17$; $s^2\approx27\,728{,}5\ne24\,687{,}5$.}
\end{ex}
